\documentclass{article}
\usepackage{graphicx}
\usepackage{amsmath}
\usepackage{amsfonts}
\usepackage{bookmark}
\usepackage{booktabs}
\usepackage{exercises}
\graphicspath{ {./images/5.3} }
\usepackage{indentfirst}
\usepackage{amsthm}
\usepackage{yhmath}
\usepackage{geometry}
\usepackage{placeins}
\usepackage{csquotes}
\usepackage{pgfplots}
\newcommand*{\oldneg}{\mathord{\sim}}
\newcommand\showdiv[1]{\overline{\smash{)}#1}}
\newcommand\Mydiv[1]{%
\smash{\raise.3ex\hbox{$\big)$}}\mkern-8mu
        \overline{\enspace\strut#1}}
\setlength{\tabcolsep}{5pt}
\theoremstyle{definition}
\newtheorem{definition}{Definition}

\setlength{\parindent}{.5in}
\setlength{\parskip}{12pt}
\renewcommand{\arraystretch}{2}
\geometry{margin=1in}
\newcommand{\?}{\stackrel{?}{=}}
\pgfplotsset{width=7cm,compat=1.18}\usepgfplotslibrary{statistics}

\title{5.4 Answers}
\begin{document}
\maketitle

\section*{1}
To divide one rational expression by another, multiply the first rational expression by the inverse of the second rational expression.

\section*{2}
You can tell that a rational expression is simplified if the numerator and the denominator don't have any other common factors than $\pm 1$.

\section*{3}
B

\section*{4}
A

\section*{5}
C

\section*{15}
\[ \frac{8x^2+10x-3}{6x^2+13x+6}=\frac{(4x-1)(2x+3)}{(3x+2)(2x+3)}=\frac{4x-1}{3x+2} \]

\section*{16}
\[ \frac{5x^2+18x-8}{10x^2-x-2}=\frac{(5x-2)(x+4)}{(5x+2)(2x-1)} \]
No simplification is possible.

\section*{17}
\[ \frac{x^3-5x^2-3x+15}{x^2-8x+15}=\frac{x^2(x)+x^2(-5)-3(x)-3(-5)}{(x-5)(x-3)}=\frac{(x^2-3)(x-5)}{(x-5)(x-3)}=\frac{x^2-3}{x-3} \]

\section*{20}
B

\section*{21-22}
The perimeter of this square is $4\times 2x=8x$, and its area is $(2x)^2=4x^2$. The perimeter of the second figure is $x+2x+x+x+3x+2x=10x$, and its area is 
$(3x\times x)+x^2=3x^2+x^2=4x^2$. 

\section*{23}
The ratio of the perimeter to the area of the first field and second field are $\frac{8x}{4x^2}=\frac{2}{x}$ and
$\frac{10x}{4x^2}=\frac{5}{2x}$, respectively. Because $\frac{5}{2x}>\frac{2}{x}$, the ratio for the second field is greater than that of the first.
I did not know which one had the lower fencing cost per unit of area.
I was also not understanding the justification for the efficiency of the package (listed on pg. 328) being the ratio of its surface area to its volume, so we should 
go over that. 

\section*{27}
\[ \frac{4(x+5)}{x^2}\times\frac{x(x+1)}{2(x+5)}=\frac{4x(x+5)(x+1)}{2x^2(x+5)}=\frac{2(x+1)}{x}=\frac{2x+2}{x} \]

\section*{30}
\[ \frac{x^2+3x-4}{x^2+4x+4}\times\frac{2x^2+4x}{x^2-4x+3}=\frac{(x^2+3x-4)(2x^2+4x)}{(x^2+4x+4)(x^2-4x+3)} \]
\[ =\frac{(x+4)(x-1)2x(x+2)}{(x+2)^2(x-3)(x-1)}=\frac{2x(x+4)}{(x+2)(x-3)} \]

\section*{33}
\[ \frac{4x^2+20x}{x^3+4x^2}\times (x^2+8x+16)=\frac{(4x^2+20x)\times (x^2+8x+16)}{x^3+4x^2} \]
\[ =\frac{4x(x+5)(x+4)^2}{x^2(x+4)}=\frac{4(x+5)(x+4)}{x} \]

\section*{34}
\[ \frac{5x^2y^3}{x^7}\div\frac{30xy^4}{y^3}=\frac{5x^2y^3}{x^7}\times\frac{y^3}{30xy^4}=\frac{5x^2y^6}{30x^8y^4}=\frac{y^2}{6x^6} \]

\section*{37}
\[ \frac{8x^2}{x+4}\div\frac{x}{2(x-4)}=\frac{8x^2}{x+4}\times\frac{2(x-4)}{x}=\frac{16x^2(x-4)}{x(x+4)}=\frac{16x(x-4)}{x+4} \]

\section*{42}
\[ \frac{x^2-8x+15}{x^2+4x}\div(x^2-x-20)=\frac{x^2-8x+15}{x^2+4x}\times\frac{1}{x^2-x-20} \]
\[ =\frac{(x-5)(x-3)}{x(x+4)(x-5)(x+4)}=\frac{x-3}{x(x+4)^2} \]

\section*{50}
(a): \[ \frac{k_1H^3V^2}{k_2H^2}=\frac{k_1HV^2}{k_2} \]


(b): \[ \frac{k_1HV^2}{k_2}=1 \] means \[ k_1HV^2=k_2 \] \[ V^2=\frac{k_2}{k_1H} \] \[ V=\sqrt{\frac{k_2}{k_1H}} \]
This means that a shorter runner has the advantage. For, the speed is calculated by taking the square root of the second constant divided by the first constant multiplied by the height. Thus, as 
the height increases, the second constant gets divided by a larger amount, meaning the speed becomes smaller.

\section*{51}
(a): The volume of a sphere is $\frac{4}{3}\pi r^3$ and the volume of a cylinder is $\pi r^2h$. We are given that $\frac{4}{3}\pi r^3=\pi r^2h$, 
meaning $h=\frac{\frac{4}{3}\pi r^3}{\pi r^2}=\frac{4}{3}r$.


(b): The surface area of the sphere is $4\pi r^2$. The surface area of the cylinder is $2\pi r(h+r)=2\pi r(\frac{4}{3}r+r)=2\pi\frac{7}{3}r^2$.


(c): The ratio of the surface area of the spherical tank to the surface area of the cylindrical tank is 
$\frac{4\pi r^2}{2\pi\frac{7}{3}r^2}=\frac{6}{7}$. The amount of material required to build the tanks is equal to the tank's surface area. Thus, because it has a lesser surface area, 
the spherical tank would take less material to build.

\section*{53}
(My answer here is off by the book's by one letter. I cannot find any errors in my answer.)\newline
We are given a cylinder with a hemisphere on each end. The length of the cylinder is $l$ and the radius of each hemisphere is $r$. Thus the two hemispheres form a sphere with radius $r$. 
The surface area of a cylinder is $2\pi r(h+r)$, and the surface area of a sphere is $4\pi r^2$. This means that the surface area of the entire shape is 
$2\pi r(l+r)+4\pi r^2=2\pi r^2+2\pi rl+4\pi r^2=6\pi r^2+2\pi rl$. The volume of the two  hemispheres is $\frac{4}{3}\pi r^3$, and the volume of the cylinder is $\pi r^2l$, meaning the 
volume of the total shape is $\frac{4}{3}\pi r^3+\pi r^2l$. Thus the ratio of the surface area to the volume of the container is:
\[ \frac{6\pi r^2+2\pi rl}{\frac{4}{3}\pi r^3+\pi r^2l}=\frac{2\pi r(3r+l)}{\frac{1}{3}\pi r^2(4r+3l)}=\frac{6(3r+l)}{r(4r+3l)} \]

\end{document}
